\documentclass[11pt]{article}
\usepackage{xltxtra}
\usepackage{bookmark}
\usepackage{hyperref}
\hypersetup{hidelinks}
\usepackage{url}
\urlstyle{tt}
\usepackage{multicol}
\usepackage{xcolor}
\usepackage{calc}
\usepackage{graphicx}
\usepackage{tikz}
\usetikzlibrary{calc}
\usepackage{fontspec}
\usepackage{xeCJK}
\usepackage{relsize}
\usepackage{xspace}
\usepackage{fontawesome5}
\usepackage{titlesec}
\usepackage{enumitem}
\usepackage{siunitx}
\usepackage{amssymb}
\usepackage{tabularx}
\usepackage{multicol}
\usepackage{fontspec}
\usepackage{fancybox}
\usepackage{float}

% ====================== 基础设置 ======================
% 取消中文字符与数字之间的间隔
\CJKsetecglue{}

% 定义美观的C++排版
\protected\def\Cpp{{C\nolinebreak[4]\hspace{-.05em}\raisebox{.28ex}{\relsize{-1}++}}\xspace}

% 全局格式设置
\setlength{\parindent}{0pt}    % 取消段落缩进
\pagenumbering{gobble}         % 取消页码显示

% 列表格式设置
\setlist[itemize]{topsep=0em, leftmargin=*}    % 项目符号列表格式
\setlist[enumerate]{topsep=0em, leftmargin=*}  % 编号列表格式

% ====================== 标题格式设置 ======================
% 一级标题格式
\titleformat{\section}
  {\LARGE\bfseries\raggedright}  % 字体：LARGE，加粗，左对齐
  {}{0em}                        % 标题前缀
  {}                             % 标题代码
  [{\color{HIT_Blue}\titlerule}]% 标题下方蓝色分隔线
\titlespacing*{\section}{0cm}{*1.2}{*1.2}  % 间距设置

% 二级标题格式  
\titleformat{\subsection}
  {\large\bfseries\raggedright}  % 字体：large，加粗，左对齐
  {}{0em}                        % 标题前缀
  {}                             % 标题代码
  []
\titlespacing*{\subsection}{0cm}{*1.2}{*1.2}  % 间距设置

% ====================== 页面布局设置 ======================
% 页面尺寸与边距
\usepackage[
    a4paper,
    left=1.2cm,
    right=1.2cm,
    top=1.5cm,
    bottom=1cm,
    nohead
]{geometry}

% 表格与正文行距设置
\renewcommand{\arraystretch}{1.2}  % 表格行距
\linespread{1.2}                   % 正文行距
\renewcommand{\CJKglue}{\hskip 0.05em}  % 中文字符间距

% ====================== 字体设置 ======================
% 字体：狮尾四季春（加糖）简体版，直接按文件从仓库 fonts/ 目录加载，无需安装到系统
% 注意：仓库自带的 SweiSpring.ttf 实为繁体 Medium 字重、SweiSpring-Bold.ttf 是损坏的空文件，
% 两者均已弃用；现用官方 github.com/max32002/swei-spring 仓库的简体 Regular 与 Bold
\setmainfont{SweiSpringSugarCJKsc-Regular.ttf}[
    Path = fonts/,
    BoldFont = SweiSpringSugarCJKsc-Bold.ttf
]
\setCJKmainfont{SweiSpringSugarCJKsc-Regular.ttf}[
    Path = fonts/,
    BoldFont = SweiSpringSugarCJKsc-Bold.ttf
]

% ====================== 主题色设置 ======================
% 哈工大校徽标准色（RGB: 0, 83, 117），参考https://xcb.hitwh.edu.cn/_upload/article/files/fa/5c/e1e0bdb340e98b473b5a0e34963c/6250efd5-f74a-47bb-b5a7-8a8b86f9451b.pdf第14页
\definecolor{HIT_Blue}{RGB}{0, 83, 117}

% ====================== 模板说明 ======================
% 这是可提交的简历模板文件，不填写真实个人信息。
% 本地私有版本请使用 main.local.tex，并确保它已加入 .gitignore。

% ====================== 基本信息 ======================
% 学院信息
\newcommand{\school}{\small{学院名称 | Faculty Name}}

% 联系方式
\newcommand{\contact}{
    \scriptsize
    \textcolor{white}{
        % 邮箱
        \faEnvelope \quad \href{mailto:your.email@example.com}{your.email@example.com}
        \hspace{4em}
        % 手机号
        \faPhone \quad  000-0000-0000
        \hspace{4em}
        % GitHub
        \faGithub \quad \href{https://github.com/yourname}{github.com/yourname}
    }
}

%%%%%%%%%%%%%%%%%%%%
% 简历正文
%%%%%%%%%%%%%%%%%%%%
\begin{document}

% ====================== 页眉 ======================
% 包含校徽、学院名称等信息
\begin{tikzpicture}[remember picture, overlay]
    \node[anchor=north, inner sep=0pt](header) at (current page.north){
        \includegraphics[width=\paperwidth]{images/header.png}
    };
    \node[anchor=west](school_logo) at (header.west){
        \hspace{0.2cm}
        \includegraphics[width=.34\textwidth]{images/hitsz_logo_name.png}
    };
    \node[anchor=east](school_name) at(header.east){
        \textcolor{white}{\textbf{\school}}
        \hspace{0.5cm}
    };
\end{tikzpicture}
\vspace{-3.5em}

% ====================== 页脚 ======================
% 包含联系方式
\begin{tikzpicture}[remember picture, overlay]
    \node[anchor=south, inner sep=0pt](footer) at (current page.south){
        \includegraphics[width=\paperwidth]{images/footer.png}
    };
    \node[anchor=center] at(footer.center){\contact};
\end{tikzpicture}

% ====================== 背景 ======================
% 添加校徽水印
\begin{tikzpicture}[remember picture, overlay]
    \node[opacity=0.05] at(current page.center){
        \includegraphics[width=0.7\paperwidth, keepaspectratio]{images/hitsz_logo.png}
    };
\end{tikzpicture}

% ====================== 个人信息 ======================
\begin{figure}[h]
    \begin{minipage}{0.82\textwidth}
        \section{\makebox[\widthof{\faAddressCard}][c]{\color{HIT_Blue}{\faAddressCard}}\quad 个人信息}
        \begin{tabularx}{\linewidth}{p{\widthof{出生日期:}}Xp{\widthof{政治面貌:}}X}
            姓\ \ \ \ \ \ \ \ 名: & 姓名 & 
            性\ \ \ \ \ \ \ \ 别: & 性别  \\
            出生年月: & YYYY年MM月 & 
            政治面貌: & 政治面貌 \\
        \end{tabularx}
    \end{minipage}
    \hspace{2em}
    \begin{minipage}{0.12\textwidth}
        \setlength{\fboxsep}{0pt}
        \doublebox{
            \begin{minipage}[c][3.2cm][c]{\linewidth}
                \centering\scriptsize 证件照
            \end{minipage}
        }
    \end{minipage}
\end{figure}
\vspace{-1em}

% ====================== 教育背景 ======================
\section{\makebox[\widthof{\faGraduationCap}][c]{\color{HIT_Blue}{\faGraduationCap}}\quad 教育背景}
\vspace{-0.5em}
\begin{table}[h!]
    \begin{tabularx}{\textwidth}{XXp{\widthof{2025年 -- 预计2029年6月毕业}}}
        学校名称 & 专业名称 & YYYY年 -- YYYY年\\
        \textbf{学分绩: XX.XX/100.00} & \textbf{学分绩排名: X/XXX} \\
    \end{tabularx}
\end{table}

% ====================== 项目经历 ======================
\section{\makebox[\widthof{\faChalkboardTeacher}][c]{\color{HIT_Blue}{\faChalkboardTeacher}}\quad 项目经历}
\vspace{0.5em}
{\large{\textbf{项目名称一}}} \hfill {技术方向 / 项目角色}\\
项目简介，说明目标、职责、关键技术、实现方法与结果。

\vspace{0.8em}
{\large{\textbf{项目名称二}}} \hfill {技术方向 / 项目角色}\\
项目简介，说明目标、职责、关键技术、实现方法与结果。

\vspace{0.8em}
{\large{\textbf{项目名称三}}} \hfill {技术方向 / 项目角色}\\
项目简介，说明目标、职责、关键技术、实现方法与结果。

% ====================== 竞赛经历 ======================
\section{\makebox[\widthof{\faTrophy}][c]{\color{HIT_Blue}{\faTrophy}}\quad 竞赛经历}
\vspace{-1em}
\begin{table}[h!]
    \begin{tabularx}{\textwidth}{Xp{\widthof{第零零负责人人人}}p{\widthof{国家级-第100名}}p{\widthof{2030年13月}}}
        \textbf{竞赛名称} & 个人/团队 & 奖项等级 & YYYY年MM月 \\
    \end{tabularx}
\end{table}

% ====================== 技能特长 ======================
\section{\makebox[\widthof{\faWrench}][c]{\color{HIT_Blue}{\faWrench}}\quad 技能特长}
\vspace{0.5em}
\begin{itemize}
    \item 编程语言：填写熟悉的编程语言、框架与工具
    \item 系统与工程：填写系统基础、工程实践、开发工具或部署经验
    \item 专业方向：填写研究方向、课程基础、项目技术栈或领域能力
    \item 数据与实验：填写数据处理、实验分析、评估指标或相关工具
    \item 综合能力：填写算法基础、协作经验、文档能力或其他优势
\end{itemize}

% ====================== 所获荣誉 ======================
\section{\makebox[\widthof{\faStar}][c]{\color{HIT_Blue}{\faStar}}\quad 所获荣誉}
\vspace{-1em}
\begin{multicols}{2}
    \begin{itemize}
        \item YYYY-YYYY学年 荣誉名称
    \end{itemize}
\end{multicols}

\end{document}
